\documentclass[11pt]{article}
\usepackage[margin=1in]{geometry}
\usepackage{natbib}
\usepackage{booktabs}
\usepackage{amsmath}
\usepackage{graphicx}
\usepackage{hyperref}
\usepackage{multirow}

\title{Time-, Tissue- and Treatment-Associated Heterogeneity in Tumour-Residing Migratory Dendritic Cells}

\author{%
  Author A$^{1,*}$, Author B$^{1}$, Author C$^{2}$, Author D$^{1,2}$ \\[6pt]
  $^{1}$Department of Immunology, University of Research, City, Country \\
  $^{2}$Institute for Cancer Research, City, Country \\
  $^{*}$Corresponding author: authora@university.edu
}

\date{}

\begin{document}
\maketitle

\begin{abstract}
CCR7$^{+}$ mature regulatory dendritic cells (mRegDCs) are recognised as key mediators of anti-tumour immunity through their capacity to migrate from tumours to draining lymph nodes (dLNs) and prime T cell responses. However, a substantial fraction of mRegDCs remains tumour-resident despite expressing migratory markers, and the phenotypic consequences of prolonged tumour dwell-time are poorly understood. Here, we employ photoconvertible mouse models to precisely timestamp and track dendritic cell migration kinetics in solid tumours. We show that mRegDCs constitute the dominant population arriving in the dLN (78--85\% of migrated DCs), yet 18\% of photoconverted DCs remain in the tumour at 72 hours post-conversion. Tumour-retained mRegDCs exhibit elevated PD-L1 (2.0-fold), reduced MHC-II (0.59-fold), and diminished co-stimulatory molecule expression relative to dLN counterparts. Single-cell RNA sequencing reveals a progressive loss of antigen-presentation and pro-inflammatory gene expression with increasing dwell-time (log$_2$FC $= -1.6$ and $-2.3$, respectively, for long-dwell cells), accompanied by intensification of regulatory programmes. Multiplexed immunohistochemistry demonstrates significant spatial co-localisation of tumour mRegDCs with PD-1$^{+}$CD8$^{+}$ T cells across murine and human tumour types (3.2--3.6-fold enrichment over random expectation). Anti-PD-L1 therapy partially reverses the regulatory phenotype, increasing MHC-II expression by 1.53-fold, boosting IL-12$^{+}$ mRegDC frequency from 12\% to 28\%, and enhancing CD8$^{+}$ T cell infiltration from 180 to 420 cells/mm$^{2}$. These findings establish that tumour dwell-time drives mRegDC functional exhaustion and that checkpoint blockade can partially restore their immunostimulatory capacity.
\end{abstract}

\section{Introduction}

Dendritic cells (DCs) are central orchestrators of adaptive anti-tumour immunity, bridging innate immune recognition with T cell activation \citep{steinman2007,mellman2011}. Among DC subsets, mature regulatory DCs (mRegDCs), characterised by co-expression of CCR7 and immunoregulatory molecules such as PD-L1, have attracted considerable attention for their dual capacity to promote T cell priming and suppress anti-tumour responses \citep{maier2020,gerhard2021}. These cells arise in the tumour microenvironment (TME) through acquisition of tumour antigens and undergo a maturation programme that includes upregulation of the chemokine receptor CCR7, which drives migration toward draining lymph nodes (dLNs) \citep{forster1999,sallusto1998}.

Classical models posit that antigen-bearing DCs rapidly transit from peripheral tissues to dLNs to present antigens and activate T cells \citep{merad2013,worbs2017}. However, emerging evidence suggests that a subset of CCR7$^{+}$ mRegDCs remains resident within the TME despite expressing migratory markers \citep{binnewies2019,veglia2021}. The functional significance of this tumour-resident population, and how prolonged exposure to the immunosuppressive TME affects mRegDC biology, has not been systematically investigated.

Several lines of evidence motivate this inquiry. First, the TME is enriched in immunosuppressive factors including IL-10, TGF-$\beta$, and prostaglandin E$_2$ that can modulate DC function \citep{gabrilovich2012,dejong2019}. Second, mRegDCs express high levels of PD-L1, positioning them at the interface between immune activation and suppression \citep{maier2020,cheng2021}. Third, the spatial relationship between tumour-resident DCs and CD8$^{+}$ T cells, which determines the potential for local antigen presentation and immune modulation, remains poorly characterised \citep{broz2014,spranger2017}.

Here, we employ photoconvertible mouse models to precisely track DC migration kinetics and dwell-time within tumours. Using a combination of flow cytometry, single-cell RNA sequencing (scRNA-seq), and multiplexed spatial analysis, we define the phenotypic and transcriptional heterogeneity of mRegDCs as a function of tissue location and tumour dwell-time. We further assess how anti-PD-L1 checkpoint immunotherapy alters tumour mRegDC biology, revealing partial reversal of the regulatory phenotype and enhanced CD8$^{+}$ T cell infiltration.

\section{Related Work}

\subsection{Dendritic cell migration and maturation}
The migration of DCs from peripheral tissues to lymph nodes is a fundamental process in adaptive immunity \citep{steinman2007,merad2013}. CCR7 and its ligands CCL19/CCL21 constitute the primary chemotactic axis driving this process \citep{forster1999,ohl2004}. Recent studies have shown that DC migration is not a binary event but rather a regulated process influenced by tissue-specific signals \citep{worbs2017}.

\subsection{mRegDCs in the tumour microenvironment}
The mRegDC programme was initially defined through single-cell transcriptomic analyses of tumour-infiltrating DCs \citep{maier2020}. These cells express a hybrid transcriptional state combining maturation markers (CCR7, FSCN1) with regulatory molecules (PD-L1, CD200) \citep{gerhard2021,cheng2021}. Subsequent work established that mRegDCs arise from both cDC1 and cDC2 precursors upon antigen uptake in the TME \citep{binnewies2019,bosteels2020}.

\subsection{Photoconvertible models for immune cell tracking}
Photoconvertible proteins such as Kaede and KikGR enable precise temporal tracking of cell migration by irreversibly shifting fluorescence from green to red upon violet light exposure \citep{tomura2008,nowotschin2009}. These systems have been applied to study DC migration in homeostatic conditions \citep{kitano2016} and during infection \citep{gerner2012}, but their application to tumour immunology has been limited.

\subsection{Checkpoint immunotherapy and DC biology}
Anti-PD-L1/PD-1 checkpoint blockade primarily targets T cell-intrinsic inhibitory signalling, but growing evidence indicates effects on DC function \citep{garris2018,oh2020}. PD-L1 on DCs can directly suppress T cell activation through both PD-1-dependent and -independent mechanisms \citep{lin2020,peng2020}. How checkpoint therapy affects the biology of tumour-resident mRegDCs specifically has not been addressed.

\section{Methods}

\subsection{Mouse models and tumour implantation}
Kaede transgenic mice (expressing the photoconvertible Kaede protein ubiquitously) were used for DC tracking experiments. MC38 colon adenocarcinoma and B16-F10 melanoma cells were implanted subcutaneously. Tumours were allowed to establish for 10 days prior to photoconversion.

\subsection{Photoconversion and DC tracking}
Tumours were surgically exposed and illuminated with 405~nm violet light to photoconvert Kaede protein from green to red fluorescence. Mice were sacrificed at 6, 24, 48, and 72 hours post-conversion, and tumours and dLNs were harvested for analysis. Photoconverted (red) DCs were tracked across tissues by flow cytometry.

\subsection{Flow cytometry}
Single-cell suspensions from tumours and dLNs were stained with antibodies against CD45, CD11c, MHC-II, CCR7, PD-L1, CD80, CD86, PD-1, and CD8. Intracellular staining for IL-12 was performed following brefeldin A treatment. mRegDCs were defined as CD11c$^{+}$MHC-II$^{+}$CCR7$^{+}$ cells. Median fluorescence intensity (MFI) values were compared between tumour and dLN compartments.

\subsection{Single-cell RNA sequencing}
Tumour mRegDCs were sorted based on photoconversion status (dwell-time bins: short $<$24~h, medium 24--48~h, long $>$48~h) and processed using the 10x Genomics Chromium platform. Data were analysed using standard pipelines including normalisation, dimensionality reduction (UMAP), clustering, and differential gene expression analysis. Gene modules for antigen presentation (H2-Aa, H2-Ab1), pro-inflammatory signalling (Il12b, Tnf), regulatory programmes (Cd274/PD-L1, Socs2), and migration (Ccr7, Fscn1) were scored.

\subsection{Multiplexed immunohistochemistry}
Formalin-fixed paraffin-embedded tumour sections from murine (MC38, B16) and human (melanoma, NSCLC) tumours were stained using multiplexed IHC panels. Spatial co-localisation of mRegDCs with PD-1$^{+}$CD8$^{+}$ T cells was quantified by measuring the fraction of mRegDCs within 30~$\mu$m of a PD-1$^{+}$CD8$^{+}$ T cell and comparing to random spatial expectation.

\subsection{Anti-PD-L1 treatment}
MC38 tumour-bearing mice were treated with anti-PD-L1 antibody or isotype control (200~$\mu$g intraperitoneally, every 3 days starting day 10 post-implantation). Tumour mRegDCs were profiled by flow cytometry at day 5 post-treatment initiation.

\subsection{Statistical analysis}
Comparisons between two groups were performed using two-tailed unpaired Student's $t$-test or Mann-Whitney $U$ test as appropriate. Spatial enrichment was assessed by permutation testing. All $p$-values less than 0.05 were considered statistically significant.

\section{Results}

\subsection{mRegDCs dominate tumour-to-dLN migration but a subset remains tumour-resident}

To track DC migration kinetics with temporal precision, we employed Kaede photoconvertible mice bearing subcutaneous tumours. Photoconversion at day 10 post-implantation allowed us to timestamp all tumour-resident cells and monitor their redistribution over time.

At 6 hours post-conversion, 85\% of photoconverted DCs remained in the tumour, with only 5\% detected in the dLN (Table~\ref{tab:migration}). By 48 hours, the tumour-resident fraction declined to 30\% while dLN-resident converted DCs peaked at 42\%. Notably, 18\% of converted DCs remained in the tumour at 72 hours, indicating a substantial tumour-resident population. Among converted DCs arriving in the dLN, mRegDCs constituted the dominant population at all time points, representing 78--85\% of migrated cells.

\begin{table}[htbp]
\centering
\caption{DC migration kinetics following tumour photoconversion. Percentage of photoconverted DCs detected in tumour and draining lymph node (dLN) at indicated time points, and proportion of mRegDCs among converted cells in the dLN.}
\label{tab:migration}
\begin{tabular}{@{}lccc@{}}
\toprule
Time post-conversion & \% Converted DCs in tumour & \% Converted DCs in dLN & \% mRegDC in dLN \\
\midrule
6 h  & 85 & 5  & 78 \\
24 h & 52 & 28 & 82 \\
48 h & 30 & 42 & 85 \\
72 h & 18 & 38 & 80 \\
\bottomrule
\end{tabular}
\end{table}

\subsection{Tumour-resident mRegDCs exhibit a regulatory phenotype distinct from dLN migrants}

To characterise the phenotypic differences between tumour-retained and migrated mRegDCs, we performed flow cytometric analysis at 48 hours post-photoconversion (Table~\ref{tab:phenotype}). Tumour mRegDCs displayed significantly elevated PD-L1 expression (MFI 6,200 vs.\ 3,100; 2.0-fold; $p < 0.001$) but markedly reduced levels of MHC-II (0.59-fold; $p < 0.001$), CD80 (0.41-fold; $p < 0.001$), CD86 (0.51-fold; $p < 0.01$), and intracellular IL-12 (0.41-fold; $p < 0.01$). CCR7 expression was also lower on tumour mRegDCs (0.62-fold; $p < 0.01$), consistent with reduced migratory capacity despite nominal CCR7 positivity.

\begin{table}[htbp]
\centering
\caption{Phenotypic comparison of tumour versus dLN mRegDCs at 48 h post-photoconversion. MFI, median fluorescence intensity.}
\label{tab:phenotype}
\begin{tabular}{@{}lcccc@{}}
\toprule
Marker (MFI) & Tumour mRegDC & dLN mRegDC & Fold difference & $p$-value \\
\midrule
CCR7           & 2,800 & 4,500 & 0.62$\times$ & $<0.01$  \\
PD-L1          & 6,200 & 3,100 & 2.0$\times$  & $<0.001$ \\
MHC-II         & 3,400 & 5,800 & 0.59$\times$ & $<0.001$ \\
CD80           & 1,200 & 2,900 & 0.41$\times$ & $<0.001$ \\
CD86           & 1,800 & 3,500 & 0.51$\times$ & $<0.01$  \\
IL-12 (intra.) & 450   & 1,100 & 0.41$\times$ & $<0.01$  \\
\bottomrule
\end{tabular}
\end{table}

\subsection{Progressive transcriptional exhaustion with increasing tumour dwell-time}

scRNA-seq analysis of tumour mRegDCs stratified by dwell-time revealed a progressive transcriptional shift (Table~\ref{tab:temporal}). Relative to freshly arrived cells ($<$24~h), medium-dwell mRegDCs (24--48~h) showed moderate reductions in antigen-presentation genes (log$_2$FC $= -0.8$) and pro-inflammatory transcripts (log$_2$FC $= -1.1$). These changes were markedly amplified in long-dwell cells ($>$48~h), with antigen-presentation genes reduced by log$_2$FC $= -1.6$ and pro-inflammatory genes by log$_2$FC $= -2.3$. Conversely, the regulatory programme (Cd274/PD-L1, Socs2) intensified progressively (log$_2$FC $= +0.5$ and $+1.2$ for medium and long dwell, respectively), while migration-associated genes declined (log$_2$FC $= -0.9$ for long dwell).

\begin{table}[htbp]
\centering
\caption{Transcriptional changes in tumour mRegDCs by dwell-time (log$_2$FC vs.\ freshly arrived, $<$24~h reference).}
\label{tab:temporal}
\begin{tabular}{@{}lccc@{}}
\toprule
Gene module & Short ($<$24 h) & Medium (24--48 h) & Long ($>$48 h) \\
\midrule
Antigen presentation (H2-Aa, H2-Ab1) & 0 (ref) & $-0.8$ & $-1.6$ \\
Pro-inflammatory (Il12b, Tnf)         & 0 (ref) & $-1.1$ & $-2.3$ \\
Regulatory (Cd274, Socs2)             & 0 (ref) & $+0.5$ & $+1.2$ \\
Migration (Ccr7, Fscn1)               & 0 (ref) & $-0.3$ & $-0.9$ \\
\bottomrule
\end{tabular}
\end{table}

\subsection{Tumour mRegDCs co-localise with PD-1$^{+}$CD8$^{+}$ T cells}

Multiplexed IHC on murine (MC38, B16) and human (melanoma, NSCLC) tumour sections revealed significant spatial co-localisation of mRegDCs with PD-1$^{+}$CD8$^{+}$ T cells (Table~\ref{tab:spatial}). Across all tumour types, 32--41\% of mRegDCs were located within 30~$\mu$m of a PD-1$^{+}$CD8$^{+}$ T cell, compared to a random expectation of 9--13\%, yielding enrichment factors of 3.2--3.6-fold. This spatial proximity was conserved across species and tumour histologies, suggesting a fundamental organisational principle of the TME.

\begin{table}[htbp]
\centering
\caption{Spatial co-localisation of mRegDCs with PD-1$^{+}$CD8$^{+}$ T cells ($<$30~$\mu$m proximity).}
\label{tab:spatial}
\begin{tabular}{@{}lccc@{}}
\toprule
Tumour type & Observed proximity & Random expectation & Enrichment fold \\
\midrule
Mouse (MC38)      & 38\% & 11\% & 3.5$\times$ \\
Mouse (B16)       & 32\% & 9\%  & 3.6$\times$ \\
Human (melanoma)  & 41\% & 13\% & 3.2$\times$ \\
Human (NSCLC)     & 35\% & 10\% & 3.5$\times$ \\
\bottomrule
\end{tabular}
\end{table}

\subsection{Anti-PD-L1 treatment reverses the regulatory phenotype of tumour mRegDCs}

Treatment of MC38 tumour-bearing mice with anti-PD-L1 antibody produced significant changes in the tumour mRegDC compartment at day 5 post-treatment (Table~\ref{tab:treatment}). mRegDC frequency among CD45$^{+}$ cells increased from 4.2\% to 6.8\% ($p < 0.01$). MHC-II expression rose by 1.53-fold (MFI 3,200 to 4,900; $p < 0.01$), and the proportion of IL-12$^{+}$ mRegDCs more than doubled from 12\% to 28\% ($p < 0.001$). Concurrently, PD-L1 MFI decreased from 6,100 to 2,800 ($p < 0.001$). These changes were accompanied by a 2.3-fold increase in CD8$^{+}$ T cell infiltration (180 to 420 cells/mm$^{2}$; $p < 0.001$).

\begin{table}[htbp]
\centering
\caption{Effect of anti-PD-L1 treatment on tumour mRegDCs and CD8$^{+}$ T cell infiltration (day 5 post-treatment).}
\label{tab:treatment}
\begin{tabular}{@{}lccc@{}}
\toprule
Parameter & Isotype control & Anti-PD-L1 & $p$-value \\
\midrule
mRegDC frequency (\% CD45$^{+}$) & 4.2  & 6.8  & $<0.01$  \\
MHC-II MFI                       & 3,200 & 4,900 & $<0.01$  \\
IL-12$^{+}$ mRegDC (\%)          & 12   & 28   & $<0.001$ \\
PD-L1 MFI                        & 6,100 & 2,800 & $<0.001$ \\
CD8$^{+}$ T cells (cells/mm$^{2}$) & 180 & 420  & $<0.001$ \\
\bottomrule
\end{tabular}
\end{table}

\section{Discussion}

Our study provides a comprehensive characterisation of mRegDC heterogeneity as a function of tissue location, tumour dwell-time, and checkpoint immunotherapy. By employing photoconvertible mouse models, we achieve temporal resolution not possible with static snapshot analyses, revealing that prolonged tumour residency drives a progressive shift from immunostimulatory to regulatory function.

The finding that 18\% of photoconverted DCs remain tumour-resident at 72 hours, despite expressing CCR7, is noteworthy. This suggests that factors within the TME actively impede mRegDC egress or that a subset of CCR7$^{+}$ cells has reduced functional migratory capacity. The lower CCR7 MFI on tumour mRegDCs (0.62-fold vs.\ dLN) supports the latter interpretation. Tumour-derived factors such as prostaglandins and adenosine have been shown to inhibit DC motility \citep{dejong2019,silva2017}, and the progressive downregulation of migration genes (Ccr7, Fscn1) with dwell-time (log$_2$FC $= -0.9$) suggests an active transcriptional programme rather than simple receptor desensitisation.

The temporal transcriptional analysis reveals a striking dwell-time-dependent exhaustion programme. Long-dwell mRegDCs lose antigen-presentation capacity (log$_2$FC $= -1.6$) while simultaneously upregulating regulatory molecules (log$_2$FC $= +1.2$). This mirrors the T cell exhaustion paradigm, in which chronic antigen exposure drives progressive functional impairment \citep{wherry2015,mclane2019}. Whether mRegDC exhaustion is reversible upon removal from the TME warrants further investigation.

The consistent spatial co-localisation of tumour mRegDCs with PD-1$^{+}$CD8$^{+}$ T cells (3.2--3.6-fold enrichment) across species and tumour types has important functional implications. Given the elevated PD-L1 on tumour mRegDCs, this proximity likely facilitates direct PD-L1/PD-1-mediated suppression of CD8$^{+}$ T cells within the TME \citep{lin2020,peng2020}. The anti-PD-L1 treatment data support this model: blocking PD-L1 reduces mRegDC PD-L1 levels and partially restores their immunostimulatory phenotype, accompanied by enhanced T cell infiltration.

A limitation of this study is the reliance on subcutaneous tumour models, which may not fully recapitulate the immunological complexity of orthotopic or spontaneous tumours \citep{olson2018}. Additionally, while we demonstrate phenotypic changes with anti-PD-L1 treatment, we did not assess whether these changes are cell-intrinsic or secondary to altered T cell activity. Future studies employing DC-specific PD-L1 deletion would clarify this distinction.

\section{Conclusion}

We demonstrate that tumour-resident mRegDCs undergo progressive functional exhaustion with increasing dwell-time, shifting from an immunostimulatory to a regulatory state characterised by elevated PD-L1, reduced antigen presentation, and diminished pro-inflammatory capacity. These cells are spatially positioned to suppress PD-1$^{+}$CD8$^{+}$ T cells within the TME. Anti-PD-L1 checkpoint therapy partially reverses this regulatory programme and enhances CD8$^{+}$ T cell infiltration. These findings identify tumour mRegDC dwell-time as a critical determinant of anti-tumour immune function and suggest that strategies targeting DC retention or rejuvenation within tumours may enhance immunotherapy efficacy.

\bibliographystyle{plainnat}
\bibliography{references}

\end{document}
